\section{器材}\label{221}

\noindent 被展示在括号中的尺寸说明公制对应值至最接近的毫米值。\\
Measurements displayed in brackets state the metric equivalent to the nearest millimetre.

\subsection{标准球台}\label{2211}

\begin{enumerate}[label=(\alph*)]
    \item 【\playingarea{}】\playingarea{}在库边\face{}内且大小应为11英尺8½英寸 × 5英尺10英寸(3569毫米 × 1778毫米)带± ½英寸(13毫米)的全部[两个]方向上的公差。\\
    【The Playing Area】The playing area is within the cushion faces and shall measure 11 ft 8½ in × 5 ft 10 in (3569 mm × 1778 mm) with a tolerance on both dimensions of +/- ½ in (13 mm).
    \item 【高度】球台从地面到库边横条上沿的高度应为2英尺10英寸(864毫米)带± ½英寸(13毫米)的公差。\\
    【Height】The height of the table from the floor to the top of the cushion rail shall be 2 ft 10 in (864 mm) with a tolerance of +/- ½ in (13 mm).
    \item 【底库和顶库】球台的两个更短的边被定义为球台的底[库](也被称为\baulk{}[库])和顶库。在带有绒毛的台呢被铺设到球台时，绒毛的光滑纹路从底库跑到顶库。\\
    【Bottom Cushion and Top Cushion】The two shorter sides of the table are defined as the Bottom (also known as Baulk) and Top Cushions of the table. Where a cloth with a nap is fitted to the table, the smooth grain of the nap runs from the Bottom Cushion to the Top Cushion.
    \item 【\baulkline{}和\baulk{}】被画在离底库\face{}29英寸(737毫米)处，且平行于它，从边库跑到边库的一条直线被称为\baulkline{}。彼直线以及[彼直线和底库]之间的空间被称为\baulk{}\footnote{虽翻译为“\baulk{}”，实则开球要在下面所说的D形区里。}。\\
    【Baulk-line and Baulk】A straight line drawn 29 in (737 mm) from the face of the Bottom Cushion, and parallel to it, running from side cushion to side cushion is called the Baulk-line. That line and the intervening space is termed Baulk.
    \item 【D形区】D形区是被标于\baulk{}内带在\baulkline{}的中点的其笔直部分的中点且带11½英寸(292毫米)半径的一个半圆。\\
    【The ``D''】The ``D'' is a semi-circle marked in Baulk with the centre of its straight section in the middle of the Baulk-line and with a radius of 11½ in (292 mm).
    \item \label{2211f}【点位】被标于D形区的每个角上的，被从\baulk{}末端观看，在右边的一个被称为黄球点且在左边的一个被称为绿球点。\\
    【Spots】Marked at each corner of the ``D'', viewed from the Baulk end, the one on the right is known as the Yellow Spot and the one on the left as the Green Spot.

    四个点位被标于球台中间\longitudinalline{}上：\\
    Four spots are marked on the centre longitudinal line of the table:
    \begin{enumerate}[label=(\roman*)]
        \item 在\baulkline{}的中点，被称为棕球点的一个；\\
        one in the middle of the Baulk-line, known as the Brown Spot;
        \item 位于顶库\face{}和底库\face{}正下方的点的正中间，被称为蓝球点的一个；\\
        one located midway between the points perpendicularly below the faces of the Top and Bottom Cushions, known as the Blue Spot;
        \item 位于蓝球点和顶库\face{}正下方的点的正中间，被称为粉球点的一个；和\\
        one located midway between the Blue Spot and a point perpendicularly below the face of the Top Cushion, known as the Pink Spot; and
        \item 离顶库\face{}正下方的点12¾英寸(324毫米)，被称为黑球点的一个。\\
        one 12¾ in (324 mm) from a point perpendicularly below the face of the Top Cushion, known as the Black Spot.
    \end{enumerate}
    \item【球袋的开设】在球台四个角的每个角处应各有一个球袋且在球台长边的中点处应每边各有一个球袋。\\
    【Pocket Openings】There shall be a pocket at each of the four corners of the table and one each at the middle of the longer sides.
\end{enumerate}

\subsection{球}

\begin{enumerate}
    \item 一套球由15颗红球，和以下彩色球：黄球、绿球、棕球、蓝球、粉球、黑球各一颗和一颗白球组成。\\
    A set of balls comprises of 15 Red balls, and one each of the following coloured balls: Yellow, Green, Brown, Blue, Pink, Black and a White. 
    \item 球应由受认证的材料构成且应各有52.5毫米带± 0.05毫米的公差的直径。\\
    The balls shall be of an approved composition and shall each have a diameter of 52.5 mm with a tolerance of +/- 0.05 mm.
    \item 若可能则球应有相等的重量但一套中最重的球和最轻的球[的重量]之间的公差应不超
    过3克。\\
    The balls shall be of equal weight where possible but the tolerance between the heaviest ball and the lightest ball in a set should be no more than 3 g.
    \item 一颗或一套球可被按球员们的共同意见或裁判的决定更换。\\
    A ball or set of balls may be changed by agreement between the players or on a decision by the referee.
\end{enumerate}

\subsection{球杆}

\noindent 球杆应长度不短于3英尺(914毫米)且应相对于带被固定在更细的一头的，被用于击打母球的一个皮头的，传统的锥形显现出无变化。\\
A cue shall be not less than 3 ft (914 mm) in length and shall show no change from the traditional tapered shape and form, with a tip, used to strike the cue-ball, secured to the thinner end.

\subsection{辅助器材}

\noindent 各种各样的架杆\footnote{包括叉头架(cross-headed rest，就是普通架杆)、大头架(butt rest)、蜘蛛架(spider rest，通常叫高架)、天鹅颈架(swan-neck rest)、加长蜘蛛架(extended spider rest)。}、加长球杆、套筒和球杆加长可被球员使用。这些可能组成在球台处\footnote{实际上平常放在球台下面。}被正常找到的器材的一部分但也包括被球员或裁判之一引入的器材。所有套筒、球杆加长和其他用于辅助击球和/或瞄准的器材都必须已得到来自相关主管机构的事先认可。\\
Various cue rests, long cues, extensions and adaptors may be used by players. These may form part of the equipment normally found at the table but also include equipment introduced by either a player or the referee. All extensions, adaptors and other devices to aid cueing and/or sighting must have received prior approval from the relevant governing body.
